% GENERATED FILE. Edit canonical lessons, then run npm run generate:books.
% canonical-source-hash: fnv1a64:21949aa8e609d4ff
% canonical-lessons: RU-C496-podguznik, RU-C496-soska, RU-C496-kolybelnaya, RU-C496-ubayukivat, RU-C496-pelenat

\chapter{A Baby at Home}
\label{ch:ru-a-baby-at-home}
\hlchaptermodality{496}

\begin{chapteropening}
\textbf{By the end of this chapter:} \emph{I can talk about looking after a baby.}

Complete the last lesson of chapter 496: I can talk about looking after a baby.
\end{chapteropening}

\section[\texorpdfstring{\ru{подгузник} (podgúznik)}{подгузник (podgúznik)}]{\ru{подгузник} (podgúznik) — a nappy, a diaper}

\label{lesson:RU-C496-podguznik}



\begin{quote}
Before the new one: say the Russian for a pram, then the Russian for to crawl.
\end{quote}

\subsection*{You'll want to know: \ru{подгузник}}
\textbf{\ru{подгузник}} (\emph{podgúznik}) — \textquotedblleft{}a nappy, a diaper\textquotedblright{}.

\begin{quote}
\ru{Папа} \ru{научился} \ru{быстро} \ru{менять} \ru{подгузники}.
\end{quote}

\begin{quote}
\emph{Pápa naúchilsya býstro menyát podgúzniki.}
\end{quote}

\begin{quote}
Dad has learnt to change nappies quickly.
\end{quote}

Masculine.

\begin{cousinweb}
\emph{\ru{подгузник}} is built on \emph{\ru{под}} 'under' and the colloquial \emph{\ru{гузно}} 'rump'; EtymDB lists the word without a source.
\end{cousinweb}

\begin{grammarlens}[title={What you've built}]
One more word for telling, explaining and describing at B1.
\end{grammarlens}

\subsection*{Your turn}
\begin{itemize}
\raggedright
  \item \emph{Say it:} \emph{podgúznik}
  \item \emph{Say it:} \emph{podgúznik}, once more
  \item \emph{Say it:} a sentence of your own with \emph{podgúznik}
  \item \emph{Recall:} say the Russian for a pram, then the Russian for to crawl, then say \emph{podgúznik} again
\end{itemize}

\begin{tcolorbox}[breakable,colback=teal!4,colframe=teal!35!black,title={Before you move on}]
What does \ru{подгузник} mean? (\textquotedblleft{}A nappy, a diaper\textquotedblright{}.) Say it once more.
\end{tcolorbox}
\section[\texorpdfstring{\ru{соска} (sóska)}{соска (sóska)}]{\ru{соска} (sóska) — a dummy, a pacifier}

\label{lesson:RU-C496-soska}



\begin{quote}
Before the new one: say the Russian for to crawl, then the Russian for a nappy.
\end{quote}

\subsection*{You'll want to know: \ru{соска}}
\textbf{\ru{соска}} (\emph{sóska}) — \textquotedblleft{}a dummy, a pacifier\textquotedblright{}.

\begin{quote}
\ru{Без} \ru{соски} \ru{малыш} \ru{не} \ru{засыпает}.
\end{quote}

\begin{quote}
\emph{Bez sóski malýsh ne zasypáyet.}
\end{quote}

\begin{quote}
The baby won't go to sleep without his dummy.
\end{quote}

Feminine. Also the teat on a baby's bottle.

\begin{cousinweb}
\emph{\ru{соска}} is built on the root of \emph{\ru{сосать}} 'to suck'; EtymDB lists \emph{\ru{соска}} without a source.
\end{cousinweb}

\begin{grammarlens}[title={What you've built}]
One more word for telling, explaining and describing at B1.
\end{grammarlens}

\subsection*{Your turn}
\begin{itemize}
\raggedright
  \item \emph{Say it:} \emph{sóska}
  \item \emph{Say it:} \emph{sóska}, once more
  \item \emph{Say it:} a sentence of your own with \emph{sóska}
  \item \emph{Recall:} say the Russian for to crawl, then the Russian for a nappy, then say \emph{sóska} again
\end{itemize}

\begin{tcolorbox}[breakable,colback=teal!4,colframe=teal!35!black,title={Before you move on}]
What does \ru{соска} mean? (\textquotedblleft{}A dummy, a pacifier\textquotedblright{}.) Say it once more.
\end{tcolorbox}
\section[\texorpdfstring{\ru{колыбельная} (kolybélnaya)}{колыбельная (kolybélnaya)}]{\ru{колыбельная} (kolybélnaya) — a lullaby}

\label{lesson:RU-C496-kolybelnaya}



\begin{quote}
Before the new one: say the Russian for a nappy, then the Russian for a dummy.
\end{quote}

\subsection*{You'll want to know: \ru{колыбельная}}
\textbf{\ru{колыбельная}} (\emph{kolybélnaya}) — \textquotedblleft{}a lullaby\textquotedblright{}.

\begin{quote}
\ru{Бабушка} \ru{поёт} \ru{внучке} \ru{колыбельную}.
\end{quote}

\begin{quote}
\emph{Bábushka poyót vnúchke kolybélnuyu.}
\end{quote}

\begin{quote}
Grandma sings her granddaughter a lullaby.
\end{quote}

Short for \emph{\ru{колыбельная} \ru{песня}}: an adjective used as a noun. Built on \emph{\ru{колыбель}} — a cradle.

\begin{cousinweb}
\emph{\ru{колыбельная}} is built on \emph{\ru{колыбель}} 'cradle', usually linked to \emph{\ru{колыхать}} 'to sway'; EtymDB has no entry.
\end{cousinweb}

\begin{grammarlens}[title={What you've built}]
One more word for telling, explaining and describing at B1.
\end{grammarlens}

\subsection*{Your turn}
\begin{itemize}
\raggedright
  \item \emph{Say it:} \emph{kolybélnaya}
  \item \emph{Say it:} \emph{kolybélnaya}, once more
  \item \emph{Say it:} a sentence of your own with \emph{kolybélnaya}
  \item \emph{Recall:} say the Russian for a nappy, then the Russian for a dummy, then say \emph{kolybélnaya} again
\end{itemize}

\begin{tcolorbox}[breakable,colback=teal!4,colframe=teal!35!black,title={Before you move on}]
What does \ru{колыбельная} mean? (\textquotedblleft{}A lullaby\textquotedblright{}.) Say it once more.
\end{tcolorbox}
\section[\texorpdfstring{\ru{убаюкивать} (ubayúkivat)}{убаюкивать (ubayúkivat)}]{\ru{убаюкивать} (ubayúkivat) — to lull to sleep}

\label{lesson:RU-C496-ubayukivat}



\begin{quote}
Before the new one: say the Russian for a dummy, then the Russian for a lullaby.
\end{quote}

\subsection*{You'll want to know: \ru{убаюкивать}}
\textbf{\ru{убаюкивать}} (\emph{ubayúkivat}) — \textquotedblleft{}to lull to sleep\textquotedblright{}.

\begin{quote}
\ru{Мама} \ru{убаюкивает} \ru{ребёнка} \ru{тихой} \ru{песней}.
\end{quote}

\begin{quote}
\emph{Máma ubayúkivayet rebyónka tíkhoy pésney.}
\end{quote}

\begin{quote}
Mum is lulling the baby to sleep with a quiet song.
\end{quote}

Imperfective; perfective \ru{убаюкать}.

\begin{cousinweb}
\emph{\ru{убаюкивать}} is built on \emph{\ru{у}-} and \emph{\ru{баюкать}} 'to lull', which is usually traced to the lullaby refrain \emph{\ru{баю}-\ru{бай}}; EtymDB lists \emph{\ru{баюкать}} without a source and has no entry for \emph{\ru{убаюкивать}}.
\end{cousinweb}

\begin{grammarlens}[title={What you've built}]
One more word for telling, explaining and describing at B1.
\end{grammarlens}

\subsection*{Your turn}
\begin{itemize}
\raggedright
  \item \emph{Say it:} \emph{ubayúkivat}
  \item \emph{Say it:} \emph{ubayúkivat}, once more
  \item \emph{Say it:} a sentence of your own with \emph{ubayúkivat}
  \item \emph{Recall:} say the Russian for a dummy, then the Russian for a lullaby, then say \emph{ubayúkivat} again
\end{itemize}

\begin{tcolorbox}[breakable,colback=teal!4,colframe=teal!35!black,title={Before you move on}]
What does \ru{убаюкивать} mean? (\textquotedblleft{}To lull to sleep\textquotedblright{}.) Say it once more.
\end{tcolorbox}
\section[\texorpdfstring{\ru{пеленать} (pelenát)}{пеленать (pelenát)}]{\ru{пеленать} (pelenát) — to swaddle}

\label{lesson:RU-C496-pelenat}



\begin{quote}
Before the new one: say the Russian for a lullaby, then the Russian for to lull to sleep.
\end{quote}

\subsection*{You'll want to know: \ru{пеленать}}
\textbf{\ru{пеленать}} (\emph{pelenát}) — \textquotedblleft{}to swaddle\textquotedblright{}.

\begin{quote}
\ru{Медсестра} \ru{показала} \ru{нам}, \ru{как} \ru{пеленать} \ru{ребёнка}.
\end{quote}

\begin{quote}
\emph{Medsestrá pokazála nam, kak pelenát rebyónka.}
\end{quote}

\begin{quote}
The nurse showed us how to swaddle the baby.
\end{quote}

Imperfective; perfective \ru{запеленать}. The cloth is a \emph{\ru{пелёнка}}.

\begin{cousinweb}
\emph{\ru{пеленать}} is built on \emph{\ru{пелена}} 'a cloth, a veil', which is inherited from Proto-Slavic \emph{pelena} 'thin skin, thin fabric' (EtymDB); EtymDB has no entry for the verb.
\end{cousinweb}

\begin{grammarlens}[title={What you've built}]
One more word for telling, explaining and describing at B1.
\end{grammarlens}

\subsection*{Your turn}
\begin{itemize}
\raggedright
  \item \emph{Say it:} \emph{pelenát}
  \item \emph{Say it:} \emph{pelenát}, once more
  \item \emph{Say it:} a sentence of your own with \emph{pelenát}
  \item \emph{Recall:} say the Russian for a lullaby, then the Russian for to lull to sleep, then say \emph{pelenát} again
\end{itemize}

\begin{tcolorbox}[breakable,colback=teal!4,colframe=teal!35!black,title={Before you move on}]
What does \ru{пеленать} mean? (\textquotedblleft{}To swaddle\textquotedblright{}.) Say it once more.
\end{tcolorbox}
